\documentclass[12pt]{article}

\usepackage[margin=1in]{geometry}
\usepackage{mathpazo}
\usepackage{parskip}
\usepackage{hyperref}

\begin{document}

\thispagestyle{empty}

\bigskip

Dear Editor,

\bigskip

We are pleased to submit our manuscript, ``Toward Self-Organized Magnonic Crystals in Ferromagnetic Shape Memory Alloys,'' for consideration in \textit{Journal of Applied Physics}.

Magnonic crystals are a core platform for wave-based computing, yet nearly all existing implementations rely on lithographic nanofabrication, which constrains them to two-dimensional geometries and to a small number of discrete states. In this paper, we investigate whether a self-assembled alternative already exists in nature. Martensitic twin variants in ferromagnetic shape memory alloys (specifically Ni-Mn-Ga) form a periodic lattice of magnetic scatterers whose periodicity is set by thermodynamics rather than lithography. This raises the question: can this naturally periodic magnetic structure function as a magnonic crystal?

We approach this question with a linearized Landau-Lifshitz-Gilbert transfer-matrix model, from which we compute the magnonic dispersion relation $\omega(q)$, the transmission spectrum, and the effect of positional disorder. Key findings include: (i)~frequency-structured transmission with strong suppression ($>40$~dB) in multiple bands, (ii)~a computed dispersion relation consistent with Bloch-wave propagation in the idealized periodic case, (iii)~dramatic changes in the transmission fingerprint under simulated magnetic field-driven variant reorientation, and (iv)~robustness of the frequency-selective attenuation against realistic positional disorder---the attenuation minima exceed 40~dB across ten independent disorder realizations at $\sigma = 10$~nm. A 2D micromagnetic simulation using mumax3 on an NVIDIA RTX 3070 Ti GPU produces results qualitatively consistent with the proposed mechanism.

The manuscript distinguishes model-supported results from experimentally testable predictions and treats potential applications as future directions.

This manuscript is appropriate for \textit{Journal of Applied Physics} because it combines magnetic materials physics, spin-wave modeling, and an experimentally testable route toward reconfigurable magnonic structures without lithographic fabrication. The paper presents theoretical results in applied magnetism that bridge materials science and magnonics, with a proposed validation pathway using Brillouin light scattering and broadband VNA-FMR on Ni-Mn-Ga thin films.

AI-assisted tools (DeepSeek and custom models) were used for manuscript drafting assistance, editorial feedback, and non-authoritative code review; all scientific content was verified by the author, who takes full responsibility for the work.

We confirm that this work is original, has not been published elsewhere, and is not under consideration by any other journal. The author has read and approved the manuscript.

We suggest the following potential referees:

\begin{itemize}
    \item Prof.\ Dirk Grundler (EPFL)---reconfigurable magnonic crystals
    \item Prof.\ Andrii Chumak (University of Vienna)---magnon spintronics
    \item Prof.\ Gaspare Varvaro (ISM-CNR, Rome)---magnetic shape memory thin films
    \item Prof.\ Paola Tiberto (INRiM, Turin)---magnetic nanostructures and magnonics
    \item Prof.\ Adekunle Adeyeye (Durham University)---nanomagnetism and magnonic crystals
\end{itemize}

Thank you for your consideration.

\bigskip

Sincerely,

\vspace{2\baselineskip}

C.\ Scott\\
Apt Systems\\
Enumclaw, Washington, USA\\
\href{mailto:cory@aptsystems.ai}{cory@aptsystems.ai}

\end{document}